%==============================================================================
% EECS 127 - Lecture 5: Linear Algebra Review - Principal Components Analysis, SVD
%==============================================================================

\documentclass[10pt, twocolumn]{article}
\usepackage[margin=0.75in, columnsep=0.3in]{geometry}

%------------------------------------------------------------------------------
% REQUIRED PACKAGES
%------------------------------------------------------------------------------
\usepackage{amsmath, amssymb, amsthm}
\usepackage{microtype}
\usepackage{enumitem}
\usepackage{fancyhdr}
\usepackage{titlesec}
\usepackage{tcolorbox}
\usepackage{booktabs}
\usepackage{tikz}
\usetikzlibrary{positioning, arrows.meta, calc}

%------------------------------------------------------------------------------
% SPACING
%------------------------------------------------------------------------------
\linespread{1.08}
\setlength{\parskip}{0.4ex plus 0.2ex minus 0.1ex}

%------------------------------------------------------------------------------
% BOX STYLES
%------------------------------------------------------------------------------
\tcbset{
    boxrule=0.8pt,
    colback=white,
    colframe=black,
    arc=0pt,
    boxsep=3pt,
    left=4pt, right=4pt, top=4pt, bottom=4pt
}

\newtcolorbox{result}{
    boxrule=0pt,
    colback=black!5,
    colframe=white,
    arc=0pt,
    boxsep=2pt,
    left=4pt, right=4pt, top=4pt, bottom=4pt
}

%------------------------------------------------------------------------------
% SECTION FORMATTING
%------------------------------------------------------------------------------
\titleformat{\section}{\large\bfseries}{\thesection.}{0.5em}{}
\titleformat{\subsection}{\normalsize\bfseries}{\thesubsection}{0.5em}{}
\titlespacing*{\section}{0pt}{1.5ex}{1ex}
\titlespacing*{\subsection}{0pt}{1ex}{0.5ex}

%------------------------------------------------------------------------------
% LIST FORMATTING
%------------------------------------------------------------------------------
\setlist{itemsep=1pt, topsep=3pt, parsep=1pt, leftmargin=1.5em}

%------------------------------------------------------------------------------
% HEADER/FOOTER
%------------------------------------------------------------------------------
\pagestyle{fancy}
\fancyhf{}
\fancyhead[L]{\small EECS 127}
\fancyhead[R]{\small Lecture 5}
\fancyfoot[C]{\small\thepage}
\renewcommand{\headrulewidth}{0.4pt}

%------------------------------------------------------------------------------
% THEOREM ENVIRONMENTS
%------------------------------------------------------------------------------
\theoremstyle{definition}
\newtheorem{definition}{Definition}
\newtheorem{example}{Example}

\theoremstyle{plain}
\newtheorem{theorem}{Theorem}
\newtheorem{lemma}{Lemma}

%------------------------------------------------------------------------------
% CUSTOM COMMANDS
%------------------------------------------------------------------------------
\newcommand{\RR}{\mathbb{R}}
\newcommand{\NN}{\mathcal{N}}
\newcommand{\Range}{\mathcal{R}}
\DeclareMathOperator{\rank}{rank}
\DeclareMathOperator{\spn}{span}
\DeclareMathOperator{\argmin}{argmin}
\DeclareMathOperator{\argmax}{argmax}
\DeclareMathOperator{\diag}{diag}

%==============================================================================

\begin{document}

\noindent
\begin{minipage}{\linewidth}
    \centering
    \textbf{\Large Lecture 5: Linear Algebra Review - Principal Components Analysis, SVD} \\[0.5em]
    \hrule
\end{minipage}
\vspace{1em}

% Textbook: Chapter 2, Sections 2.5-2.6

%==============================================================================
\section{Principal Component Analysis}
%==============================================================================

In many applications, data lives in very high-dimensional spaces (thousands or millions of dimensions). \textbf{Principal Component Analysis} (PCA) is a technique to find the most important lower-dimensional structure in the data. Formally, given data in $\RR^d$, we want to find an underlying $p$-dimensional linear structure where $p \ll d$.

%------------------------------------------------------------------------------
\subsection{Data Matrix and Covariance Matrix}
%------------------------------------------------------------------------------

Suppose we have data points $\vec{x}_1, \ldots, \vec{x}_n \in \RR^d$. We assume the data is \textbf{zero-mean}: $\frac{1}{n} \sum_{i=1}^{n} \vec{x}_i = \vec{0}$.

We organize these into a \textbf{data matrix} $X$ where data points form the rows:
\begin{result}
\[
    X \triangleq \begin{pmatrix} \vec{x}_1^\top \\ \vdots \\ \vec{x}_n^\top \end{pmatrix} \!\in \RR^{n \times d}, \quad X^\top = \begin{pmatrix} \vec{x}_1 & \!\cdots\! & \vec{x}_n \end{pmatrix} \!\in \RR^{d \times n}
\]
\end{result}

We define the \textbf{covariance matrix} $C \in \RR^{d \times d}$ by
\[
    C \triangleq \frac{1}{n} X^\top X = \frac{1}{n} \sum_{i=1}^{n} \vec{x}_i \vec{x}_i^\top
\]

Since $X^\top X$ is symmetric, we have $C \in \mathbb{S}^d$ (the set of $d \times d$ symmetric matrices). By Lecture~4, $C$ has real eigenvalues and orthogonal eigenvectors.

%------------------------------------------------------------------------------
\subsection{The PCA Optimization Problem}
%------------------------------------------------------------------------------

We want to find a unit vector $\vec{w}_1 \in \RR^d$ (with $\|\vec{w}_1\|_2 = 1$) such that projecting the data onto $\vec{w}_1$ preserves as much structure as possible.

\textbf{Step 1: Write the projection error.}

The projection of $\vec{x}_i$ onto $\spn(\vec{w}_1)$ is $(\vec{w}_1^\top \vec{x}_i)\vec{w}_1$. The squared error for a single point is:
\[
    e_i^2 = \|\vec{x}_i - (\vec{w}_1^\top \vec{x}_i)\vec{w}_1\|_2^2
\]

The average error across all data points is:
\[
    \text{err}(\vec{w}_1) = \frac{1}{n} \sum_{i=1}^{n} \|\vec{x}_i - (\vec{w}_1^\top \vec{x}_i)\vec{w}_1\|_2^2
\]

\textbf{Step 2: Expand the error.}

Since $\|\vec{w}_1\|_2 = 1$, expanding the squared norm gives:
\begin{align*}
    e_i^2 &= \|\vec{x}_i\|_2^2 - 2(\vec{x}_i^\top \vec{w}_1)^2 + (\vec{x}_i^\top \vec{w}_1)^2 \|\vec{w}_1\|_2^2 \\
    &= \|\vec{x}_i\|_2^2 - (\vec{x}_i^\top \vec{w}_1)^2
\end{align*}

So the average error becomes:
\[
    \text{err}(\vec{w}_1) = \frac{1}{n} \sum_{i=1}^{n} \|\vec{x}_i\|_2^2 - \frac{1}{n} \sum_{i=1}^{n} (\vec{x}_i^\top \vec{w}_1)^2
\]

\textbf{Step 3: Convert minimization to maximization.}

The first term $\frac{1}{n}\sum \|\vec{x}_i\|_2^2$ is a constant (it does not depend on $\vec{w}_1$). So minimizing the error is equivalent to:
\begin{result}
\[
    \max_{\substack{\vec{w}_1 \in \RR^d \\ \|\vec{w}_1\|_2 = 1}} \frac{1}{n} \sum_{i=1}^{n} (\vec{x}_i^\top \vec{w}_1)^2
\]
\end{result}

\begin{tcolorbox}[colframe=black!50]
\textbf{Intuition:} Minimizing reconstruction error is equivalent to maximizing the variance of the projections. The best direction $\vec{w}_1$ is the one along which the data is most spread out.
\end{tcolorbox}

%------------------------------------------------------------------------------
\subsection{Reduction to the Rayleigh Quotient}
%------------------------------------------------------------------------------

\textbf{Step 4: Rewrite in matrix form.}

We can factor the objective using the covariance matrix:
\begin{align*}
    \frac{1}{n} \sum_{i=1}^{n} (\vec{x}_i^\top \vec{w}_1)^2 &= \frac{1}{n} \sum_{i=1}^{n} \vec{w}_1^\top \vec{x}_i \vec{x}_i^\top \vec{w}_1 \\
    &= \vec{w}_1^\top \underbrace{\left(\frac{1}{n} \sum_{i=1}^{n} \vec{x}_i \vec{x}_i^\top\right)}_{C} \vec{w}_1 \\
    &= \vec{w}_1^\top C \vec{w}_1
\end{align*}

The quantity $\vec{w}_1^\top C \vec{w}_1$ (for $\|\vec{w}_1\|_2 = 1$) is called the \textbf{Rayleigh quotient} of $C$ at $\vec{w}_1$.

\textbf{Step 5: Solve via the spectral theorem.}

Since $C \in \mathbb{S}^d$, by the spectral theorem we can write $C = U \Lambda U^\top$ where $U$ is orthogonal and $\Lambda = \diag(\lambda_1, \ldots, \lambda_d)$ with $\lambda_1 \ge \cdots \ge \lambda_d$.

Substituting $\vec{y} = U^\top \vec{w}_1$ (note $\|\vec{y}\|_2 = \|\vec{w}_1\|_2 = 1$ since $U$ preserves norms):
\[
    \vec{w}_1^\top C \vec{w}_1 = \vec{y}^\top \Lambda \vec{y} = \sum_{i=1}^{d} \lambda_i y_i^2
\]

Since $\sum y_i^2 = 1$ and $y_i^2 \ge 0$, this is a convex combination of the eigenvalues. Therefore:

\begin{result}
\[
    \lambda_{\min}(C) \le \vec{w}^\top C \vec{w} \le \lambda_{\max}(C) \quad \text{for all } \|\vec{w}\|_2 = 1
\]
\end{result}

The upper bound is achieved when $\vec{y} = \vec{e}_1$ (all weight on the largest eigenvalue), i.e., when $\vec{w}_1 = U\vec{e}_1$ is the eigenvector corresponding to $\lambda_{\max}(C)$.

%------------------------------------------------------------------------------
\subsection{The First Principal Component}
%------------------------------------------------------------------------------

\begin{tcolorbox}
\textbf{Theorem} (First Principal Component):

The first principal component $\vec{w}_1^\star$ is the eigenvector of $C = \frac{1}{n}X^\top X$ corresponding to its largest eigenvalue $\lambda_{\max}(C)$. The minimum reconstruction error is:
\[
    \min_{\|\vec{w}_1\|_2 = 1} \text{err}(\vec{w}_1) = \frac{1}{n} \sum_{i=1}^{n} \|\vec{x}_i\|_2^2 - \lambda_{\max}(C)
\]
\end{tcolorbox}

\begin{tcolorbox}[colframe=black!50]
\textbf{Higher principal components.} The second principal component $\vec{w}_2$ is found by repeating the same optimization, but restricting $\vec{w}_2 \perp \vec{w}_1$. This yields the eigenvector for the second-largest eigenvalue, and so on. The SVD (next section) provides a clean way to compute all principal components simultaneously.
\end{tcolorbox}

%==============================================================================
\section{Singular Value Decomposition}
%==============================================================================

The \textbf{Singular Value Decomposition} (SVD) generalizes the eigendecomposition to \emph{any} matrix, not just symmetric or square ones. It is one of the most important tools in linear algebra and data science.

%------------------------------------------------------------------------------
\subsection{Definition}
%------------------------------------------------------------------------------

\begin{tcolorbox}
\textbf{Definition} (SVD):

Let $A \in \RR^{m \times n}$ have rank $r$. A \textbf{singular value decomposition} of $A$ is:
\[
    A = U \Sigma V^\top
\]
where:
\begin{itemize}
    \item $U \in \RR^{m \times m}$ is orthogonal, with columns $\vec{u}_1, \ldots, \vec{u}_m$ called \textbf{left singular vectors}
    \item $V \in \RR^{n \times n}$ is orthogonal, with columns $\vec{v}_1, \ldots, \vec{v}_n$ called \textbf{right singular vectors}
    \item $\Sigma \in \RR^{m \times n}$ has entries $\sigma_1 \ge \cdots \ge \sigma_r > 0$ on the diagonal and zeros elsewhere, where $\sigma_i$ are the \textbf{singular values}
\end{itemize}
\end{tcolorbox}

%------------------------------------------------------------------------------
\subsection{Three Forms of the SVD}
%------------------------------------------------------------------------------

The SVD can be written in three equivalent ways, each useful in different contexts.

\textbf{Full SVD.} The definition above: $A = U\Sigma V^\top$ with $U \in \RR^{m \times m}$, $\Sigma \in \RR^{m \times n}$, $V \in \RR^{n \times n}$.

If $A$ is \textbf{tall} ($m > n$) with full column rank $n$:
\[
    A = U \begin{pmatrix} \Sigma_n \\ 0_{(m-n) \times n} \end{pmatrix} V^\top
\]

If $A$ is \textbf{wide} ($m < n$) with full row rank $m$:
\[
    A = U \begin{pmatrix} \Sigma_m & 0_{m \times (n-m)} \end{pmatrix} V^\top
\]

\textbf{Compact SVD.} Keep only the $r$ nonzero singular values and their corresponding singular vectors:

\begin{result}
\[
    A = U_r \Sigma_r V_r^\top
\]
\end{result}

where $U_r \in \RR^{m \times r}$, $\Sigma_r = \diag(\sigma_1, \ldots, \sigma_r) \in \RR^{r \times r}$, $V_r \in \RR^{n \times r}$.

\textbf{Dyadic SVD.} Express $A$ as a sum of rank-1 matrices (dyads):

\begin{result}
\[
    A = \sum_{i=1}^{r} \sigma_i \vec{u}_i \vec{v}_i^\top
\]
\end{result}

Each term $\sigma_i \vec{u}_i \vec{v}_i^\top$ is a rank-1 matrix called a \textbf{dyad}. The dyads are ordered by importance (largest $\sigma_i$ first).

%------------------------------------------------------------------------------
\subsection{Constructing the SVD}
%------------------------------------------------------------------------------

How do we actually find $U$, $\Sigma$, $V$? We use the eigendecomposition of $A^\top A$.

\textbf{Step 1: Eigendecompose $A^\top A$.}

The matrix $A^\top A \in \RR^{n \times n}$ is symmetric and positive semidefinite. It has rank $r$ and thus $r$ positive eigenvalues. Order them:
\[
    \lambda_1 \ge \cdots \ge \lambda_r > \lambda_{r+1} = \cdots = \lambda_n = 0
\]
with corresponding orthonormal eigenvectors $\vec{v}_1, \ldots, \vec{v}_n$. These are the \textbf{right singular vectors}.

\textbf{Step 2: Compute singular values.}

Define the singular values as the square roots of the eigenvalues:
\begin{result}
\[
    \sigma_i \triangleq \sqrt{\lambda_i} > 0 \quad \text{for } i = 1, \ldots, r
\]
\end{result}

\textbf{Step 3: Compute left singular vectors.}

For $i = 1, \ldots, r$, define:
\[
    \vec{u}_i \triangleq \frac{A\vec{v}_i}{\sigma_i}
\]

This gives $r$ orthonormal vectors in $\RR^m$. To fill out $U \in \RR^{m \times m}$, extend $\{\vec{u}_1, \ldots, \vec{u}_r\}$ to an orthonormal basis for $\RR^m$ using Gram-Schmidt (the remaining $m - r$ vectors span $\NN(A^\top)$).

\begin{theorem}
In the SVD construction above, $\{\vec{u}_1, \ldots, \vec{u}_m\}$ is an orthonormal set and $A = U\Sigma V^\top$.
\end{theorem}

\begin{tcolorbox}[colframe=black!50]
\textbf{Non-uniqueness.} The SVD is not unique. The Gram-Schmidt extension is arbitrary, repeated singular values allow different eigenvector choices, and even distinct eigenvalues only determine eigenvectors up to a sign $\vec{v} \mapsto -\vec{v}$.
\end{tcolorbox}

%------------------------------------------------------------------------------
\subsection{Geometry of the SVD}
%------------------------------------------------------------------------------

The key geometric insight is that multiplying by $A$ is equivalent to three operations:
\begin{enumerate}
    \item \textbf{Rotate/reflect} by $V^\top$: maps the standard basis to the $V$-basis
    \item \textbf{Scale} by $\Sigma$: stretches each axis by $\sigma_i$
    \item \textbf{Rotate/reflect} by $U$: aligns the scaled axes in the output space
\end{enumerate}

\begin{tcolorbox}[colframe=black!50]
\textbf{Geometric picture:} $V^\top$ maps the unit circle to the unit circle (rotation). $\Sigma$ maps the unit circle to an axis-aligned ellipse with semi-axes $\sigma_1, \ldots, \sigma_r$. $U$ rotates this ellipse to its final orientation.
\end{tcolorbox}

\textbf{Action on a general vector.} To understand $A\vec{x}$, write $\vec{x}$ in the $V$-basis: $\vec{x} = \alpha_1 \vec{v}_1 + \cdots + \alpha_n \vec{v}_n$. Then by linearity:
\[
    A\vec{x} = \alpha_1 \sigma_1 \vec{u}_1 + \cdots + \alpha_r \sigma_r \vec{u}_r
\]

Components beyond $r$ are annihilated (they are in $\NN(A)$).

\textbf{Maximum and minimum scaling.} Since $\sigma_1 \ge \cdots \ge \sigma_r > 0$:
\begin{result}
\[
    \max_{\|\vec{x}\|_2 = 1} \|A\vec{x}\|_2 = \sigma_1, \qquad \min_{\substack{\|\vec{x}\|_2 = 1 \\ \vec{x} \notin \NN(A)}} \|A\vec{x}\|_2 = \sigma_r
\]
\end{result}

The largest singular value $\sigma_1$ is the maximum stretching factor of $A$, and $\sigma_r$ is the minimum nonzero stretching factor. If $A$ is not full column rank, some nonzero vectors are sent to $\vec{0}$.

%------------------------------------------------------------------------------
\subsection{Connection to PCA}
%------------------------------------------------------------------------------

The SVD directly computes PCA. If $X \in \RR^{n \times d}$ is our zero-mean data matrix and we compute the SVD $X = U\Sigma V^\top$, then:

\begin{itemize}
    \item The covariance matrix is $C = \frac{1}{n}X^\top X = \frac{1}{n}V\Sigma^\top \Sigma V^\top = V \left(\frac{\Sigma^\top \Sigma}{n}\right) V^\top$
    \item The columns of $V$ are the eigenvectors of $C$, i.e., the \textbf{principal components}
    \item The eigenvalues of $C$ are $\lambda_i = \sigma_i^2 / n$
    \item The first principal component is the column of $V$ corresponding to the largest singular value $\sigma_1$
\end{itemize}

\begin{tcolorbox}[colframe=black!50]
\textbf{In practice:} We compute PCA by taking the SVD of $X$ (not by forming $X^\top X$ and eigendecomposing it). The SVD is numerically more stable and efficient.
\end{tcolorbox}

%==============================================================================

\vfill

\begin{center}
\rule{0.5\linewidth}{0.4pt}

\textit{Key Takeaways}
\end{center}

\begin{enumerate}
    \item \textbf{PCA finds directions of maximum variance:} minimizing reconstruction error is equivalent to maximizing $\vec{w}^\top C \vec{w}$ subject to $\|\vec{w}\|_2 = 1$.
    \item \textbf{Rayleigh quotient bound:} for $A \in \mathbb{S}^d$ and $\|\vec{w}\|_2 = 1$, we have $\lambda_{\min} \le \vec{w}^\top A \vec{w} \le \lambda_{\max}$, with extrema achieved by eigenvectors.
    \item \textbf{First principal component} is the eigenvector of $C = \frac{1}{n}X^\top X$ for its largest eigenvalue.
    \item \textbf{SVD decomposes any matrix:} $A = U\Sigma V^\top$ (full), $A = U_r \Sigma_r V_r^\top$ (compact), $A = \sum_{i=1}^{r} \sigma_i \vec{u}_i \vec{v}_i^\top$ (dyadic).
    \item \textbf{SVD construction:} eigendecompose $A^\top A$ to get $V$ and $\sigma_i = \sqrt{\lambda_i}$, then $\vec{u}_i = A\vec{v}_i / \sigma_i$.
    \item \textbf{Geometry:} the SVD decomposes $A$ into rotation--scaling--rotation; $\sigma_1$ is the maximum stretching factor.
    \item \textbf{SVD computes PCA:} the columns of $V$ from the SVD of $X$ are the principal components.
\end{enumerate}

\end{document}
